% ============================================================
%  cap06/6.2_recomendaciones.tex
% ============================================================
\section{Recomendaciones}
A partir del análisis forense de los hallazgos operativos y del despliegue integral del sistema de misión crítica \textit{Democra}, se establecen las siguientes directrices, proyecciones arquitectónicas y recomendaciones tecnológicas para los futuros ciclos de iteración, evolución e investigación científica sobre la plataforma. Estas recomendaciones trascienden las propuestas de mejora convencionales, apuntando a la inserción del sistema en paradigmas de computación distribuida de vanguardia y estándares de resiliencia algorítmica a prueba de los vectores de amenaza del mañana.

\begin{enumerate}
    \item \textbf{Evolución hacia el Patrón Event Sourcing Puro mediante Apache Kafka:}
    Si bien la aproximación actual basada en la captura de eventos a nivel de triggers en formato \textit{JSONB} ha demostrado suficiencia para la auditoría y la reconstrucción lógica del estado, se recomienda encarecidamente transitar la capa de persistencia transaccional hacia un modelo de \textit{Event Sourcing} puro, orquestado a través de un sistema de mensajería distribuido y log estructurado, tal como Apache Kafka o Redpanda. En esta arquitectura, el estado actual de cualquier entidad (el \textit{Materialized View}) sería estrictamente derivado de una secuencia inmutable de eventos de dominio anexados de forma \textit{append-only}. Esta migración erradicaría definitivamente cualquier riesgo de pérdida de historial por mutaciones en la base de datos relacional y habilitaría un desacoplamiento asintótico total entre la ingestión de datos a altísima velocidad (con escrituras secuenciales en el disco de $\mathcal{O}(1)$) y los procesos de lectura, los cuales se podrían consumir y proyectar en estructuras NoSQL heterogéneas, optimizadas según los patrones de acceso (Command Query Responsibility Segregation o CQRS). Dicha transición dotará al sistema de un mecanismo de recuperación ante desastres (Disaster Recovery) trivializado a la mera reproducción \textit{replay} del \textit{topic} de eventos desde el \textit{offset} cero.

    \item \textbf{Implementación Rigurosa de Prácticas de Chaos Engineering:}
    Las pruebas de carga tradicionales no logran exponer los modos de fallo sistémico latentes en arquitecturas distribuidas modernas. Por consiguiente, se insta vehementemente a abandonar las simulaciones estáticas de concurrencia y a instaurar una disciplina continua de \textit{Chaos Engineering}, inspirada en los principios del \textit{Simian Army}. Este marco de validación empírica requerirá la inyección programática y estocástica de fallos en el clúster de producción (por ejemplo: interrupción artificial de nodos, degradación controlada de la latencia en enlaces de red mediante \textit{traffic shaping}, saturación sintética de memoria, y terminación aleatoria de réplicas de bases de datos). El objetivo matemático de esta disciplina es confirmar la resiliencia inherente de \textit{Democra} y garantizar que los algoritmos de \textit{failover}, los patrones \textit{Circuit Breaker}, y las políticas de reintento con \textit{Exponential Backoff} actúen con una precisión determinística, impidiendo las cascadas de fallos catastróficas y manteniendo la disponibilidad continua de los servicios esenciales a pesar de las hecatombes de la infraestructura subyacente.

    \item \textbf{Adopción de Criptografía Post-Cuántica (PQC) para la Integridad a Largo Plazo:}
    Considerando que \textit{Democra} almacena datos de carácter civil, clínico y financiero de extrema sensibilidad, se debe prever que la amenaza teórica que representan las computadoras cuánticas a gran escala (el escenario del \textit{Store Now, Decrypt Later} mediante la aplicación del algoritmo de Shor) se materialice en la próxima década. Por tanto, se aconseja proyectar la migración progresiva de la infraestructura de clave pública (PKI) interna, los túneles de cifrado de transporte (TLS) y, fundamentalmente, la firma criptográfica del árbol de registros inmutables, hacia esquemas de encriptación que sean resistentes a ataques cuánticos. La integración de los estándares del NIST, tales como los algoritmos de encapsulación de claves basados en celosías matemáticas (Lattice-based cryptography como CRYSTALS-Kyber) y los esquemas de firma digital CRYSTALS-Dilithium o SPHINCS+, dotará a la plataforma de un nivel de inmunidad de grado militar contra adversarios de computación cuántica, resguardando el secreto perfecto hacia adelante (\textit{Perfect Forward Secrecy}) y garantizando que las trazas de auditoría de los movimientos gubernamentales sigan siendo innegables para la posteridad.

    \item \textbf{Orquestación Analítica Avanzada e Integración con Redes Neuronales Profundas:}
    Se debe impulsar el acoplamiento de los repositorios de datos masivos consolidados (Data Lakes) hacia arquitecturas de inteligencia artificial integradas en el flujo de valor. En lugar de un mero módulo de Inteligencia de Negocios (BI) bidimensional, se propone la construcción de canales de datos (pipelines) directos hacia ecosistemas de cómputo en clúster (como Apache Spark) para el entrenamiento estocástico de modelos predictivos y redes neuronales profundas (Deep Learning) aplicadas. Estos tensores matemáticos podrían analizar series de tiempo de alta dimensionalidad sobre los patrones de donación, métricas de engagement del voluntariado y el procesamiento de lenguaje natural de los casos clínicos reportados. Al desplegar arquitecturas MLOps para la inferencia en tiempo real, \textit{Democra} transcendería de una plataforma de gestión reactiva hacia un oráculo predictivo capaz de prescribir, de forma completamente autónoma, reubicaciones tácticas del personal voluntario y optimizaciones dinámicas del capital financiero antes de que se presenten déficits operativos sistémicos.

    \item \textbf{Aislamiento Absoluto de Micro-inquilinos Vía Computación Confidencial (Confidential Computing):}
    A medida que la aplicación ingrese en los ámbitos corporativos hiper-regulados y agencias de inteligencia, el modelo RLS (Row Level Security) requerirá un salto evolutivo hacia el hardware. Se recomienda la exploración y adaptación del motor transaccional al paradigma de la \textit{Confidential Computing}, aprovechando enclaves seguros basados en procesadores avanzados (tales como Intel SGX o AMD SEV-SNP). Bajo esta topología, la información altamente confidencial procesada por Democra se mantendrá encriptada no solo en tránsito (in-transit) y en reposo (at-rest), sino también durante su procesamiento en memoria RAM (in-use). Este grado absoluto de opacidad criptográfica impediría que incluso el administrador del hipervisor de la nube pública o los administradores de los sistemas subyacentes pudiesen inspeccionar o exfiltrar los datos descifrados o las claves de cifrado en el tiempo de ejecución, cerrando definitivamente las puertas a ataques de inyección a nivel kernel, elevación de privilegios en el sistema host, e inspección profunda de la memoria.
\end{enumerate}
